\chapter{Fundamentação Teórica}

\section{Bulas de Medicamentos no Brasil}

\subsection{Estrutura e Exigências da ANVISA}

\subsection{Protocolos, Resoluções e Orientações Oficiais}

% [?] pode haver seção sobre acessibilidade ao leigo ou comparação bula/profissional vs. bula/paciente

\section{Modelos de Linguagem}

\subsection{Arquitetura Transformer}

\subsection{Modelos de Grande e Pequeno Porte}

\subsection{Português Brasileiro e Terminologia Médico-Farmacêutica}

\subsection{Tokenização e Representação de Termos Especializados}

\subsection{Execução Local, Quantização e Eficiência Computacional}

\subsection{Privacidade e Soberania Tecnológica}

% [?] fine-tuning, destilação e instruction tuning entram apenas se forem usados ou delimitados como trabalhos futuros

\section{Geração Aumentada por Recuperação (RAG)}

\subsection{Pipeline RAG}

\subsection{Representação Vetorial e Indexação}

% [?] pode haver seção sobre estratégias de chunking ou re-ranking se forem abordadas na metodologia

\section{Avaliação de Sistemas de Pergunta e Resposta}

\subsection{Factualidade, Grounding, Segurança e Utilidade}

\subsection{Métricas Léxicas}

\subsection{Métricas Semânticas}

\subsection{Métricas de RAG}

\subsection{Avaliação Humana e Taxonomias de Erro}
